\chapter{Neurological Exam}
\section*{What Is This Test?} A neurological exam assesses brain, spinal cord, nerve, and muscle function.\section*{Alternative Names} Neurologic examination; neurological assessment.\section*{Principle} Mental status, cranial nerves, strength, sensation, reflexes, coordination, and gait localize dysfunction.\section*{Why This Test Is Ordered} It evaluates weakness, numbness, headache, seizure, dizziness, confusion, tremor, or balance change.\section*{Primary vs Secondary Test} Examination is primary; imaging, blood, cerebrospinal fluid, or electrodiagnostic tests are selected secondarily.\section*{How This Test Is Performed} The clinician observes speech and movement and tests eye movements, power, sensation, reflexes, coordination, and walking.\section*{What Preparation Is Needed} None; bring symptom timing, medicine, and prior neurologic records.\section*{Understanding Results} Findings are interpreted by pattern and urgency; sudden focal deficits require emergency assessment.\section*{Follow-up Tests} CT, MRI, EEG, lumbar puncture, EMG, or targeted laboratory tests may follow.\section*{References} \begin{enumerate}\item MedlinePlus. Neurological Exam.\item American Academy of Neurology. Neurologic examination resources.\end{enumerate}\clearpage\section*{Understanding Results and Follow-up}
The laboratory report should identify the specimen, method, units, reference interval or decision limit, and whether the result is positive, negative, abnormal, or inconclusive. Reference intervals vary with age, sex, pregnancy, specimen type, assay, and laboratory; a result outside a range is not automatically a diagnosis, and a result within a range does not always exclude disease. Discuss unexpected results with the ordering clinician, who can decide whether repeat or confirmatory testing, treatment, or referral is appropriate. Do not change medicines or delay urgent care based on an isolated result.
\section*{Regulatory Status and Authorization Date}This chapter describes a test category, not a specific commercial product. FDA, CDSCO, EU IVDR, and MHRA approval or registration dates therefore cannot be assigned without the assay manufacturer, product name, instrument, and intended use. For laboratory-developed tests, an individual agency approval date may not exist. See the Preface for the interpretation of regulatory status.
\clearpage
